\section{Related Work}
\label{sec:related_work}

\textbf{Weight Ensembling and Task Arithmetic.} 
Directly merging the parameters of specialized neural networks fine-tuned from a shared initialization allows running a single forward pass instead of execution of multiple models, saving massive computation. Simple linear interpolation of weights \cite{wortsman2022model} and Task Arithmetic \cite{ilharco2022editing} (adding task vectors linearly) are widely used but suffer from inter-task interference and representation misalignment. Advanced static schemes like TIES-Merging \cite{yadav2023ties} (resolving sign conflicts) and DARE \cite{yu2023language} (random delta pruning) mitigate interference but remain inherently static and non-adaptive, computing a single set of merged weights that cannot adapt dynamically to varying task compositions at test-time.

\textbf{Dynamic Routing and Mixture-of-Experts.} 
To enable sample-specific parameter ensembling, test-time dynamic model merging and routing frameworks have emerged. Operating at the model or adapter level, several methods optimize routing coefficients at test-time over a calibration split, such as FoldMerge \cite{coordinatewarping} (using coordinate warping) and QWS-Merge \cite{qws2025} (quantum-inspired superposition). However, recent work exposes severe vulnerabilities in these paradigms: (i) \emph{Layer-Averaging Collapse} \cite{l3collapse} shows that layer-wise routing is an over-parameterized redundancy; (ii) \emph{Vectorization Collapse} \cite{priorclassical} shows that standard dynamic routers collapse when transitioning to batch-independent streams ($B=1$) due to the lack of batch-smoothing, necessitating complex regularization; and (iii) the \emph{Dynamic Routing Paradox} shows that parametric routers optimized on tiny calibration sets ($N=64$) inevitably overfit. Parameter-Free Subspace Routing (PFSR) \cite{pfsr2026} sidesteps optimization collapse by using classification weights as anchors, but its reliance on unweighted cosine similarity assumes a flat, isotropic space, which is geometrically and information-theoretically misspecified.

\textbf{Information Geometry in Deep Learning.} 
Information geometry treats parameter spaces as Riemannian manifolds with local metrics defined by the Fisher Information Matrix (FIM) \cite{amari1998natural}, which measures expected local curvature of the Kullback-Leibler divergence. The diagonal empirical Fisher Information Matrix (dFIM) has been widely adopted in deep learning, notably in Elastic Weight Consolidation (EWC) \cite{kirkpatrick2017overcoming} to prevent catastrophic forgetting. While prior works use Fisher Information for offline, static weight ensembling (e.g., Fisher Merging \cite{matena2021merging}), our work is the first to utilize Fisher Information dynamically at test-time as a coordinate-warping metric tensor for parameter-free subspace routing, unifying information-geometric curvature with dynamic test-time model merging.
